\documentclass[tikz,border=12pt,12pt]{standalone}
\usepackage{ctex}
\usepackage{xcolor}
\usepackage{tikz}
\usetikzlibrary{positioning, calc, arrows.meta, shadows, shapes.geometric, shapes.misc, fit, backgrounds, matrix}

% Color palette
\definecolor{layer1}{RGB}{219,234,254}
\definecolor{layer1d}{RGB}{59,130,246}
\definecolor{layer2}{RGB}{220,252,231}
\definecolor{layer2d}{RGB}{34,197,94}
\definecolor{layer3}{RGB}{254,243,199}
\definecolor{layer3d}{RGB}{234,179,8}
\definecolor{layer4}{RGB}{252,231,243}
\definecolor{layer4d}{RGB}{236,72,153}
\definecolor{layer5}{RGB}{237,233,254}
\definecolor{layer5d}{RGB}{139,92,246}
\definecolor{layer6}{RGB}{255,237,213}
\definecolor{layer6d}{RGB}{249,115,22}
\definecolor{accent}{RGB}{37,99,235}
\definecolor{accentdeep}{RGB}{29,78,216}
\definecolor{bg}{RGB}{251,252,253}
\definecolor{textmuted}{RGB}{91,100,117}
\definecolor{borderc}{RGB}{226,232,240}

\begin{document}
\begin{tikzpicture}[
    node distance=0.9cm,
    font=\sffamily,
    every node/.style={inner sep=0pt, outer sep=0pt},
    layerbox/.style={
        rectangle, rounded corners=8pt,
        minimum width=14cm, minimum height=1.3cm,
        draw=#1, thick,
        fill=#1!25,
        drop shadow={shadow xshift=2pt, shadow yshift=-2pt, fill=gray!25},
        align=left,
    },
    layertitle/.style={
        font=\bfseries\large,
        text=#1,
        yshift=6pt,
    },
    layersub/.style={
        font=\small,
        text=gray!70,
        yshift=-10pt,
    },
    arrowdown/.style={
        -{Stealth[scale=1.2]},
        thick,
        draw=gray!50,
    },
]

% Background
\node [rectangle, fill=bg, minimum width=17cm, minimum height=21cm, rounded corners=12pt, draw=borderc, very thick] at (0,0) {};

% Title
\node [font=\bfseries\Huge, text=accentdeep] at (0, 9.2) {Cinematography IR};
\node [font=\large, text=textmuted] at (0, 8.5) {电影摄影分层中间表示 · 架构信息图};
\draw [thick, draw=accent!60, rounded corners=2pt] (-4.5, 8.1) -- (4.5, 8.1);

% ========== Layer 1: Narrative Intent ==========
\node [layerbox=layer1d, minimum height=1.5cm] (l1) at (0, 6.8) {};
\node [layertitle=blue!80, anchor=west] at ($(l1.west)+(10pt,0)$) {\textcircled{1} \enspace 叙事意图层 \quad Narrative Intent};
\node [layersub, anchor=west] at ($(l1.west)+(10pt,0)$) {戏剧问题 · 情绪节拍 · 主题 · 观众状态 · Emotional Arc};

% Arrow 1
\draw [arrowdown] (l1.south) -- ++(0, -0.55);

% ========== Layer 2: Staging & Blocking ==========
\node [layerbox=layer2d, minimum height=1.5cm] (l2) at (0, 4.95) {};
\node [layertitle=green!80, anchor=west] at ($(l2.west)+(10pt,0)$) {\textcircled{2} \enspace 场面调度层 \quad Staging \& Blocking};
\node [layersub, anchor=west] at ($(l2.west)+(10pt,0)$) {人物/道具 · 空间标记 · 视线 · 动作轨迹 · 灯光布置 · Beat};

% Arrow 2
\draw [arrowdown] (l2.south) -- ++(0, -0.55);

% ========== Layer 3: Shot Design ==========
\node [layerbox=layer3d, minimum height=1.5cm] (l3) at (0, 3.1) {};
\node [layertitle=yellow!80!orange, anchor=west] at ($(l3.west)+(10pt,0)$) {\textcircled{3} \enspace 镜头设计层 \quad Shot Design};
\node [layersub, anchor=west] at ($(l3.west)+(10pt,0)$) {景别 · 机位角度 · 构图策略 · Coverage 角色 · 叙事目的 · 深度分层};

% Arrow 3
\draw [arrowdown] (l3.south) -- ++(0, -0.55);

% ========== Layer 4: Camera Program ==========
\node [layerbox=layer4d, minimum height=1.5cm] (l4) at (0, 1.25) {};
\node [layertitle=magenta!80, anchor=west] at ($(l4.west)+(10pt,0)$) {\textcircled{4} \enspace 相机程序层 \quad Camera Program};
\node [layersub, anchor=west] at ($(l4.west)+(10pt,0)$) {位姿 · 焦距/光圈 · 对焦 · 曝光 · 18种运镜 · 缓动曲线 · 灯光引用};

% Arrow 4
\draw [arrowdown] (l4.south) -- ++(0, -0.55);

% ========== Layer 5: Continuity ==========
\node [layerbox=layer5d, minimum height=1.5cm] (l5) at (0, -0.6) {};
\node [layertitle=violet!80, anchor=west] at ($(l5.west)+(10pt,0)$) {\textcircled{5} \enspace 连续性约束层 \quad Continuity Constraints};
\node [layersub, anchor=west] at ($(l5.west)+(10pt,0)$) {180°轴线 · 银幕方向 · 互反视线 · 30°规则 · 桥接例外 · 空间重置};

% Arrow 5
\draw [arrowdown] (l5.south) -- ++(0, -0.55);

% ========== Layer 6: Backend ==========
\node [layerbox=layer6d, minimum height=1.2cm] (l6) at (0, -2.2) {};
\node [layertitle=orange!80, anchor=west] at ($(l6.west)+(10pt,0)$) {\textcircled{6} \enspace 渲染适配层 \quad Renderer / DCC / VP Adapter};
\node [layersub, anchor=west] at ($(l6.west)+(10pt,0)$) {Blender · Unreal · Three.js · USD · 虚拟制片};

% ========== Side panel: IR Core label ==========
\begin{scope}[on background layer]
\node [rectangle, fill=accent!10, minimum width=0.7cm, minimum height=9.2cm, rounded corners=4pt]
  at ($(l1.east)+(0.8cm, -2.9cm)$) (corelabel) {};
\node [rotate=90, font=\bfseries\small, text=accentdeep] at (corelabel) {IR CORE · 模型 + 验证 + 分析};
\end{scope}

% ========== Bottom section: modules & CLI ==========
\draw [thick, draw=gray!30, dashed] (-7.2, -3.2) -- (7.2, -3.2);

\node [font=\bfseries\Large, text=accentdeep] at (0, -4.0) {核心模块};

% Module cards row
\tikzset{modulecard/.style={
    rectangle, rounded corners=8pt,
    minimum width=3.3cm, minimum height=1.6cm,
    draw=borderc, thick,
    fill=white,
    drop shadow={shadow xshift=1.5pt, shadow yshift=-1.5pt, fill=gray!20},
    align=center,
}}

\node [modulecard] (m1) at (-5.2, -5.3) {
    \textbf{\textcolor{blue!70}{model.rs}}\\
    \small \textcolor{gray!70}{IR 数据模型}\\
    \small \textcolor{gray!70}{~880 行·25+ 类型}
};
\node [modulecard] (m2) at (-1.7, -5.3) {
    \textbf{\textcolor{green!70}{validate.rs}}\\
    \small \textcolor{gray!70}{结构语义验证}\\
    \small \textcolor{gray!70}{~990 行·40+ 规则}
};
\node [modulecard] (m3) at (1.7, -5.3) {
    \textbf{\textcolor{violet!70}{analyze.rs}}\\
    \small \textcolor{gray!70}{连续性分析}\\
    \small \textcolor{gray!70}{~215 行·4 条规则}
};
\node [modulecard] (m4) at (5.2, -5.3) {
    \textbf{\textcolor{orange!70}{diagnostic.rs}}\\
    \small \textcolor{gray!70}{诊断系统}\\
    \small \textcolor{gray!70}{~95 行·JSON 可序列化}
};

% Second row
\node [modulecard] (m5) at (-5.2, -7.2) {
    \textbf{\textcolor{teal!70}{io.rs}}\\
    \small \textcolor{gray!70}{YAML / JSON 加载}\\
    \small \textcolor{gray!70}{自动格式分派}
};
\node [modulecard] (m6) at (-1.7, -7.2) {
    \textbf{\textcolor{red!70}{report.rs}}\\
    \small \textcolor{gray!70}{人类可读摘要}\\
    \small \textcolor{gray!70}{镜头清单输出}
};
\node [modulecard] (m7) at (1.7, -7.2) {
    \textbf{\textcolor{purple!70}{main.rs}}\\
    \small \textcolor{gray!70}{cine-ir CLI}\\
    \small \textcolor{gray!70}{5 个子命令}
};
\node [modulecard] (m8) at (5.2, -7.2) {
    \textbf{\textcolor{cyan!70}{docs/}}\\
    \small \textcolor{gray!70}{中文摄影教材}\\
    \small \textcolor{gray!70}{14 章 mdBook}
};

% Connect layers to module cards
\draw [dotted, gray!40, thick] (l6.south) -- ++(0, -0.4) -| (m1.north);

% ========== CLI section ==========
\draw [thick, draw=gray!30, dashed] (-7.2, -8.2) -- (7.2, -8.2);

\node [font=\bfseries\Large, text=accentdeep] at (0, -8.9) {cine-ir CLI};

\tikzset{clichip/.style={
    rectangle, rounded corners=6pt,
    minimum width=2.5cm, minimum height=0.9cm,
    draw=accent!40,
    fill=accent!10,
    font=\small\ttfamily,
    text=accentdeep,
    align=center,
}}

\node [clichip] (c1) at (-5.0, -9.9) {validate};
\node [clichip] (c2) at (-2.5, -9.9) {analyze};
\node [clichip] (c3) at (0.0, -9.9) {inspect};
\node [clichip] (c4) at (2.5, -9.9) {normalize};
\node [clichip] (c5) at (5.0, -9.9) {schema};

\node [font=\small, text=textmuted] at (0, -10.7) {结构验证 · 连续性分析 · 摘要检视 · 格式归一化 · JSON Schema};

% Footer
\draw [thick, draw=gray!20] (-7.2, -11.3) -- (7.2, -11.3);
\node [font=\small\ttfamily, text=gray!50] at (0, -11.75) {cinematography-ir v0.1.0 \quad|\quad Rust 2021 \quad|\quad MIT};

\end{tikzpicture}
\end{document}
